% !TEX program = pdflatex
% Integration proposal v2 (post-GPT-round2) for DM-deficient UDG stress
% test into the ECT preprint.
%
% Changes vs v1:
%   - Dropped Proposal D (footnote to internal working PDF).
%   - Weakened accent on G_eff(X) as "main rescue path"; G_eff is listed
%     equal-footing with history-dependent Xi and orientation-sector.
%   - Removed "isolated-field DM-deficient dwarfs should exist" as a
%     preprint prediction; replaced with honest statement that the theory
%     does not yet make such a prediction.
%   - Executive recommendation now reads A + B only.
%
% This is still internal working document, NOT for publication.

\documentclass[11pt]{article}
\usepackage[margin=2.2cm]{geometry}
\usepackage{amsmath,amssymb}
\usepackage{booktabs}
\usepackage{hyperref}
\usepackage{xcolor}
\hypersetup{colorlinks=true,linkcolor=blue,urlcolor=blue}

\title{Proposal (v2): DM-deficient UDG stress test --- minimal integration\\
into the ECT preprint}
\author{Internal working document --- NOT for publication}
\date{2026-04-20}

\begin{document}
\maketitle

\section*{Executive recommendation}

\textbf{Recommendation:} Integrate at minimal-commitment level:
\begin{enumerate}
\item One new row in {\tt tab:future\_discriminants} (ND-13).
\item One short paragraph (~150 words) in the existing galactic-closure
discussion, referencing the ND-13 entry.
\item \textbf{No standalone subsection.} Claiming a derivation would
overstate the current theoretical state.
\item \textbf{No footnote reference} to any internal working notes as
semi-official supplement to the preprint.
\end{enumerate}

\textbf{Rationale.} DF4/FCC~224 observations are public, well-known,
and will be raised by any careful reviewer regardless of whether the
preprint addresses them. Pre-emptive honest disclosure with a clear
``open discriminant'' label is safer than silent omission. A full
subsection would invite scrutiny of unsupported claims. A short honest
acknowledgement is defensible and strengthens the paper's scientific
credibility. What is \emph{not} recommended is any claim that the
current preprint explains DF4, or that the dynamical $G_{\rm eff}(\phi)$
degree of freedom already provides an intra-preprint rescue mechanism:
the preprint does not derive the connection between $G_{\rm eff}(\phi)$
and the practical galactic closure.

\bigskip\hrule\bigskip

\section{Text proposed for integration}

\subsection{Proposal A: {\tt tab:future\_discriminants} row ND-13}

\textit{Goal:} Add one row to the discriminants table already in the
preprint.

\textbf{Row format} (matching the existing table structure):

\begin{quote}\small
\textbf{ND-13: Dark-matter-deficient ultra-diffuse galaxies.}
\textit{Class: galactic kinematics stress test.}

\textbf{Observable.} Stellar (and GC) velocity dispersion
$\sigma_{\rm obs}$ of ultra-diffuse galaxies identified via top-heavy
GCLF and low $M_{\rm dyn}/M_{\rm bar}$ signatures. Reference sample:
DF4 (\cite{Danieli2022DF4}, $\sigma=6.3^{+2.5}_{-1.6}$ km\,s$^{-1}$,
$M_\star\simeq 1.5\times 10^{8}M_\odot$, $R_e\simeq 1.6$~kpc);
FCC~224 (\cite{Buzzo2025FCC224,Tang2025FCC224},
$\sigma\simeq 7.8$ km\,s$^{-1}$); DF2
(\cite{Trujillo2019DF2,Shen2021DF2TRGB}) which remains observationally
ambiguous. DF4, and probably FCC~224, exhibit tracer kinematics
consistent with baryon-only dynamical mass in the probed region; DF2
remains observationally ambiguous.

\textbf{Current closure prediction.} The practical Level-B galactic
closure~(eq.~\ref{eq:ect_rotation_closure_main}) at the baseline
$g^\dagger_{\rm eff}=g^\dagger_0$ tends to the MOND-like
$\sigma\sim\sqrt{g_N g^\dagger_0}\,R_{1/2}$ branch for these systems,
giving $\sigma\sim 18$--$25$ km\,s$^{-1}$ (Wolf proxy within the
S\'ersic half-mass radius). Baryon-only Newtonian values lie at
$\sigma\sim 8$--$14$ km\,s$^{-1}$, much closer to the observed
$\sigma\sim 6$--$8$ km\,s$^{-1}$. These quoted ranges are diagnostic
proxy-level comparisons (Wolf/Jeans-style stress-test estimates), not
yet production-level ECT fits with matched tracer-family, aperture,
and full anisotropic modelling.

\textbf{Status.} Open stress test of the current Level-B closure, not
a solved consequence of it. The minimal density-only Level-B
ansatz~(eq.~\ref{eq:gdagger_env_main}) for $\Xi$ does not naturally
accommodate both the low-dispersion systems (DF4, probably FCC~224)
and matched DM-rich controls such as NGC~5846-UDG1
(\cite{Forbes2021NGC5846UDG1}, $\sigma=17\pm 2$ km\,s$^{-1}$ at
comparable $M_\star$, $R_e$, and group density) within a single
simple environmental scalar law. The presently known sample is small
and not selection-neutral: several candidates were identified through
unusual GC-population markers, so the class properties should be
treated as provisional rather than population-level conclusions.

\textbf{Possible theory-level directions for a future completion.}
Dynamical $G_{\rm eff}(\phi)$ in eq.~\ref{eq:gen_EE},
history-dependent extensions of $\Xi$ within the broader
ECT-compatible class noted in \S\ref{sec:env_sector}, and
orientation-sector coupling through $\alpha_{\rm ECT}$ in
eq.~\ref{eq:gen_EE}. These belong to structures already present in the
broader framework, but none is presently derived for the galactic
closure; addressing them is an open problem linked to OP3.

\textbf{Distinctive prediction.} A future rescue mechanism should make
a definite statement about whether isolated-field DM-deficient dwarfs
are expected or forbidden in ECT; at present the theory does not yet
make that prediction. Additional observational tightening is expected
from matched DM-rich UDG controls (NGC~5846-UDG1, VCC~1287) and from
isolated-field searches.

\textbf{Comparative standing.} $\Lambda$CDM with bullet-dwarf
formation reproduces the observed $\sigma$
values~\cite{Ogiya2018DF2tidal,Moreno2022DMfree,vanDokkum2022bullet}
through partial dark-halo stripping; any ECT-compatible rescue path
will need to reach comparable explanatory power to remain competitive
on this class of objects.
\end{quote}

\subsection{Proposal B: short paragraph in body of paper}

\textit{Location:} After the SPARC rotation-curve fit discussion and
before the discussion of $\Xi$ environmental modulation.

\textit{Proposed text} (~150 words):

\begin{quote}\small
The matched-mass population of pressure-supported ultra-diffuse
galaxies contains both DM-rich objects (NGC~5846-UDG1 with
$\sigma=17\pm 2$ km\,s$^{-1}$ at $M_\star\sim 10^{8}M_\odot$ and
$R_e\sim 2$~kpc, consistent with the RAR) and DM-deficient objects
(DF4, FCC~224, with $\sigma\sim 6$--$8$ km\,s$^{-1}$ at comparable
$M_\star$ and $R_e$). The current practical closure of the $\phi$-first
galactic sector~(eq.~\ref{eq:ect_rotation_closure_main}) is a Level-B
interpolation with Newtonian and deep-MOND asymptotics; the exact
$\Xi$-law remains open (OP3). At the baseline
$g^\dagger_{\rm eff}=g^\dagger_0$, the closure is consistent with the
matched DM-rich control at the proxy-comparison level, but tends to
over-predict $\sigma$ on the DM-deficient case by a factor of order
three. DF4 and probably FCC~224 therefore constitute empirical
stress tests of the present Level-B closure: they are not a solved
consequence of current ECT, but an open discriminator for any future
completion of $\Xi(\rho_{\rm env},\bar\phi_{\rm env},\ldots)$ and/or
of the broader time-dependent/orientation sector. We list this as
entry ND-13 in tab.~\ref{tab:future_discriminants}. The $\Lambda$CDM
bullet-dwarf interpretation remains a viable comparator.
\end{quote}

\subsection{Proposal C: no standalone subsection}

A standalone subsection titled, e.g., ``DM-deficient UDGs in ECT''
would suggest that ECT provides a derived explanation. It does not.
Including such a subsection without a derivation would be
scientifically incorrect at the current state of the theory.
\textbf{Recommendation: do not include.}

\subsection{(Dropped) Proposal D: footnote to internal working notes}

Dropped. Internal working documents should not be referenced from the
published preprint as semi-official supplements. If they become
relevant in a future version of the theory, they will be cited as
proper supplementary material at that time.

\section{What the integration delivers to the paper}

\begin{enumerate}
\item \textbf{Honest disclosure} of a known observational tension, which
any astronomy-literate reviewer will test the paper against.
\item \textbf{Scientific credibility} from explicit acknowledgement of
where the theory currently stands, rather than silent omission.
\item \textbf{Concrete falsifiable discriminant}: the ND-13 entry gives
future observers a clear target for constraining ECT.
\item \textbf{Pointers, not claims}, to candidate theory-level
directions (dynamical $G_{\rm eff}(\phi)$, history-dependent $\Xi$,
orientation-sector), explicitly flagged as not-yet-derived.
\item \textbf{Fair comparison to $\Lambda$CDM}: explicit acknowledgement
that bullet-dwarf dark-matter-stripped simulations exist and are
competitive on this class of objects.
\end{enumerate}

\section{What the integration does \emph{not} claim}

\begin{enumerate}
\item Does not claim that ECT explains DF4 or FCC~224.
\item Does not claim that a derived rescue mechanism exists.
\item Does not introduce any new phenomenology outside the preprint.
\item Does not advertise $G_{\rm eff}(\phi)$ as a solved intra-preprint
rescue path; it is listed equal-footing with other theory-level
directions, all labeled as not-yet-derived.
\item Does not reference internal working notes as a semi-official
supplement.
\item Does not claim that ``$\bar\phi_{\rm local}\to 0^-$'' implies
Newton-like DF4 within the current closure.
\end{enumerate}

\section{Word count estimate}

\begin{tabular}{lr}
\toprule
Element & Words added to preprint\\
\midrule
{\tt tab:future\_discriminants} ND-13 row & $\sim 330$\\
Body paragraph (Proposal B) & $\sim 160$\\
\midrule
\textbf{Total (A + B)} & $\sim 490$\\
\bottomrule
\end{tabular}

Current preprint $\sim 670$ pages; addition $\sim 0.1\%$ in volume,
scientifically valuable, low risk.

\section{What happens if we do not integrate}

\textbf{Risk 1:} arXiv endorser sees the DF4/FCC~224 literature, expects
the preprint to acknowledge it, does not find anything, and declines
endorsement with a note such as ``should at least acknowledge known
stress tests.''

\textbf{Risk 2:} A first informed reader of the posted preprint writes
a public comment citing DF4/FCC~224. The paper loses the opportunity
to frame the issue on its own terms.

\textbf{Risk 3:} The claimed galactic-sector success rings hollow to
experts aware of specific dwarfs that strain the closure.

All three risks are mitigated by the minimal Proposal A + B integration.
None is mitigated by waiting or by hoping reviewers will not notice.

\section{What happens if the integration is accepted}

\textbf{Gain 1:} The paper passes the honesty test for astronomers
reviewing the galactic-sector claims.

\textbf{Gain 2:} The paper provides a concrete target for future work,
anchoring the programmatic research agenda on ECT's galactic sector
to a specific OP3-linked open problem.

\textbf{Gain 3:} The arXiv endorser sees a preprint that handles its
own challenges openly, which typically improves endorsement prospects.

\textbf{Gain 4:} Any future ECT publication that resolves
$\Xi(\rho_{\rm env},\bar\phi_{\rm env},\ldots)$ (or provides a
derivation of a suppression route via $G_{\rm eff}(\phi)$,
non-quasi-static $\phi$-dynamics, or orientation-sector effects) can
cite ND-13 as a previously flagged open discriminator rather than as
a newly-discovered problem.

\section{Decision to be made by Valera}

\begin{enumerate}
\item[Decision 1:] Integrate Proposal A (table entry ND-13)?
\emph{Recommended: yes.}

\item[Decision 2:] Integrate Proposal B (short body paragraph)?
\emph{Recommended: yes, if A is accepted.}

\item[Decision 3:] Integrate Proposal C (standalone subsection)?
\emph{Recommended: no.}
\end{enumerate}

\textbf{Recommended minimum:} A + B.

Once Valera confirms which proposals to accept, the exact LaTeX insert
locations (anchors in {\tt ECT\_preprint.tex}) and final wording can
be finalized and committed via the standard backup--edit--compile--push
workflow. GPT review of the resulting inserts should precede the push.

\end{document}
